\chapter{绪论}

\section{研究背景与意义}

在量化投资研究中，股票收益的形成机制具有显著的多因素驱动特征。无论是价值、成长、质量、动量、波动率、流动性等经典风格因子，还是宏观环境、行业景气度、公司事件与市场情绪等扩展信息，都会在不同程度上影响股票的定价与未来收益\textsuperscript{[13][14][15]}。随着数据采集技术、量化平台和研究数据库的发展，研究者已经不再局限于少量核心指标，而是可以同时获取财务报表、日频 K 线、技术指标、行业分类、指数成分变化、分红公告、重大事项乃至外部宏观变量等多源信息。由此带来的直接结果是：股票研究对象正在从“少量因子解释收益”逐步走向“高维因子联合刻画收益”，而高维、多源、长时间窗口已经成为当前量化研究的重要特征。

在这一背景下，股票因子数据的复杂性主要体现在三个层面。第一，横截面维度上，同一时点往往需要同时跟踪数十只、数百只甚至更多股票，不同股票之间还存在行业、规模、风格与联动关系。第二，因子维度上，单只股票在一个交易时点通常对应多个基本面、技术面和扩展信息因子，因子之间往往高度相关，容易产生冗余和共线性。第三，时间维度上，股票因子并非静态不变，而会随着政策环境、市场流动性、宏观周期与情绪状态变化而发生阶段性漂移\textsuperscript{[10][14][17]}。因此，股票因子研究的基础对象已经不再是单纯的二维表格，而是天然包含“股票—因子—时间”三重结构的高维数据系统。若忽略这种结构，仅把它视为普通矩阵来处理，往往很难完整描述市场中真实存在的多维耦合关系\textsuperscript{[6][9][15]}。

传统多因子研究中，更常见的做法是先将原始数据展平为“样本$\times$特征”的矩阵，再使用主成分分析、逐步回归、相关系数筛选、聚类或打分模型进行降维与选股\textsuperscript{[13][14]}。这类方法具有计算简单、工程实现成熟、结果便于解释等优点，因此长期以来一直是实务和学术研究中的主流路径。然而，其隐含前提通常是：样本之间相对独立、特征之间的关系较稳定、时间维度只作为附加标签存在。对于真实股票因子问题而言，这些前提并不总是成立。首先，同一类因子的有效性可能在牛熊转换、货币宽松或风险偏好变化阶段出现明显差异；其次，不同行业和不同样本池中的股票，对相同因子的敏感度并不一致；再次，当因子数量持续扩张时，仅依靠二维矩阵压缩容易把横截面差异、因子共振关系和时间状态变化混合在一起，从而损失原始结构信息，并削弱模型的解释性和稳定性\textsuperscript{[6][7][9][15]}。

高维因子问题不仅带来建模层面的困难，也直接带来统计与计算层面的挑战。一方面，因子数量增加后，多重共线性、噪声累积、过拟合风险以及参数不稳定等问题会更加突出，高维因子模型研究也因此强调“如何在保留主要信息的前提下控制潜在维度”\textsuperscript{[15][16]}。另一方面，若仍使用单一二维框架处理多维数据，则很容易在矩阵化过程中破坏原始结构，甚至导致降维结果虽然数值上可用，但经济含义变弱。Lettau 在高维因子模型研究中指出，高维张量化数据可以在大幅压缩维度的同时保留主要变异，并且潜在成分往往具有较清晰的经济解释\textsuperscript{[15]}。这说明在金融场景中，降维问题的关键不只是“把维数降下来”，还包括“能否在降维后仍然保留结构和解释”。

张量分解方法为解决这一矛盾提供了一条更自然的建模路径。张量是矩阵的高阶推广，可以直接在原始多模态结构下开展表示学习，而不必在分析起点就把所有维度粗糙展平\textsuperscript{[4][6][7]}。在经典张量方法中，CP 分解通过若干秩一张量之和表达原始张量，适合提取共享的潜在成分；Tucker 分解通过核张量与多个因子矩阵的乘积刻画不同模态之间更复杂的交互作用，具有更强的灵活性。已有研究表明，张量方法不仅适用于图像识别、数据压缩和结构保持降维，也已经逐步扩展到金融时序预测、股票波动分析、市场可预测性研究和股票相关结构识别等任务\textsuperscript{[8][9][10][11][12][17]}。对于股票因子研究来说，张量分解的意义在于，它既可以完成高维因子压缩，又能够同时输出股票模态、因子模态和时间模态的潜在表示，从而为股票联动模式识别、因子共振关系发现和时间状态切换分析提供统一基础。

对中国 A 股市场而言，这种方法价值尤为明显。A 股样本本身具有指数成分动态变化、停牌与复牌并存、财务披露存在时滞、市场风格轮动明显、行业结构差异较大等特点。如果仍然依赖静态、二维和单窗口式的分析框架，就很容易把样本池边界变化、时间漂移和多因子共振这些问题割裂开来处理。相比之下，以三阶张量形式组织股票因子数据，能够更自然地把样本边界、因子集合与时间窗口纳入同一表达体系，也更有利于后续进行跨样本池比较与滚动窗口稳健性分析。换言之，张量方法不仅是一个新的“算法工具”，更对应了一种更适合 A 股复杂数据结构的研究组织方式。

从理论意义上看，本文将张量分解应用于股票因子降维与模式发现，一方面能够在金融数据分析语境下进一步检验高阶张量建模的适用性，比较 CP、Tucker 与 PCA 三种路径在结构保持、信息提取和结果解释上的差异边界\textsuperscript{[4][9][15][16]}；另一方面，也能够把张量分解从信号处理、图像分析等领域更明确地引入到股票因子研究中，为“高维多因子体系如何在保留结构信息前提下实现有效压缩”提供更系统的方法论支持。特别是当研究目标不再局限于单纯预测，而同时关注因子解释、股票结构和时间状态时，张量方法所具有的多模态表示能力就更具理论价值。

从实践意义上看，若能在高维股票因子集合中识别出少量信息密度更高、结构解释更清晰的低维核心表示，就可以降低选股模型复杂度，减轻多重共线性问题，增强风格解释与风险监控的稳定性\textsuperscript{[14][15]}。进一步地，若张量分解输出的时间模态能够识别市场状态切换、风格轮动和异常阶段的结构突变，那么它不仅有助于提升单次实验结果的可读性，也有助于构建更稳健、更可复用的量化研究流程\textsuperscript{[10][11][12][17]}。因此，围绕 A 股正式样本开展股票因子张量建模、降维与模式发现研究，不仅具有方法上的创新空间，也具有较强的现实应用价值。

\section{国内外研究现状}

张量分解的数学思想可以追溯到 Hitchcock 对多元数组表示形式的研究。Hitchcock 在 1927 年给出了将高阶张量表示为多个乘积项求和的基本形式，为后续规范分解奠定了数学起点\textsuperscript{[1]}。随后，Tucker 在三模因子分析中提出核张量结构，使高阶数据可以通过不同模态的潜在因子矩阵进行表达\textsuperscript{[2]}；Harshman 又从多模态因子分析角度发展了 PARAFAC 方法，使 CP 分解逐步成为高阶结构分析中的经典模型\textsuperscript{[3]}。Kolda 与 Bader 的综述系统总结了张量分解的数学定义、算法框架和应用场景，被广泛视为该领域的重要基础文献\textsuperscript{[4]}。胡胜龙关于张量分解唯一性的研究则进一步梳理了 CP 分解的可辨识条件和理论边界，为高维张量模型的解释性提供了更扎实的理论支撑\textsuperscript{[5]}。

从方法综述角度看，夏虹对张量方法及其应用的综述说明，张量建模的最大优势在于能够直接处理多维结构数据，而不需要在建模初期将多模态信息粗糙展平\textsuperscript{[6]}。Rabanser、Shchur 与 Günnemann 的综述则从机器学习视角系统介绍了 CP、Tucker 及其在主题模型、关系数据和时序问题中的应用，指出张量分解已从纯数学工具演化为高维数据分析中的通用方法族\textsuperscript{[7]}。这些研究共同表明，张量方法的核心价值并不仅仅在于“高维”，更在于“结构保持”。

在金融数据分析方向，国内已有研究开始尝试将张量模型引入具体金融任务，而不再满足于把高维数据简单展平成二维矩阵。岳欣的研究以停牌情形下的股票波动分析为切入点，选取 2018 年 6 月 1 日至 2020 年 12 月 31 日上证 50 成分股中多个停牌股票的多项交易指标，将其组织为张量形式，并采用基于 CP 分解的 CP-WOPT 算法对缺失值进行联合填充\textsuperscript{[8]}。该研究解决的是“停牌导致样本缺失后如何保留多指标联合结构”的问题，说明张量分解不仅适用于降维，也可用于缺失恢复与后续波动分析；但其任务重心仍是停牌补全和波动刻画，并未进一步讨论降维后的因子排序或结构解释。曾亚丽则把研究对象放在证券市场指数预测上，先利用单个证券市场多元时间序列构造二阶张量，再利用多个证券市场相关技术指标构造三阶张量，并分别引入$(2D)^2$PCA+PCA、MPCA 和 ENW-MPCA 等特征提取算法\textsuperscript{[10]}。其中，$(2D)^2$PCA+PCA 试图解决二维张量特征提取后维数仍然偏高的问题，MPCA 用于解决向量化方法无法保持三维内部结构的问题，而 ENW-MPCA 则进一步针对不同特征向量对预测贡献程度不同这一现象，通过特征值归一化加权提升预测精度。该研究较为系统地说明了张量表示在时间序列预测中的优势，但它服务的是市场指数预测任务，而不是股票因子降维与横截面解释。

从更一般的统计学习角度看，马少沛等关注的是高阶张量如何在保持结构的前提下降维\textsuperscript{[9]}。该文把 Zhong 等提出的张量充分降维思路从二阶推广到更一般的高阶情形，并以三阶张量为例提出结构转换算法和结构保持算法，重点解决高维张量降维中协方差矩阵奇异、计算复杂度过高以及向量化操作破坏原始结构的问题。作者还将所提方法与 K-means、t-SNE、多维主成分分析、多维判别分析和张量 SIR 等方法进行比较，发现结构保持算法在分类精度上更具优势。这一研究虽然应用场景是图像识别，但它提供了一个对本文非常重要的启示：高维张量分析的关键不只是“降维”，更是“在降维过程中保留原有结构关系”。此外，Brandi 等尝试利用张量化分解研究股票相关结构，强调股票之间的联动关系不应只由二维收益相关系数刻画；Spelta 利用张量分解与 links forecast 讨论金融市场可预测性；Wang 等则从高维向量自回归时序建模出发，说明张量分解可用于复杂时序系统的参数压缩与潜在结构提取\textsuperscript{[11][12][17]}。这些研究共同说明，张量方法在金融场景中的价值已经从单一的数学建模扩展到缺失恢复、结构识别、预测和参数压缩等多个方面。

在多因子选股领域，现有中文研究更多聚焦于因子筛选、赋权和宏观修正，而不是显式利用张量结构。黄宏运等从股票选取中的多因子问题出发，收集了 2937 只股票的 10 项影响因子，既包括市净率、市盈率、资产负债比率等基本面指标，也包括 ADR、RSI、BIAS、OBV、RSV 等技术面指标，并以相对收益率为被解释变量建立多元线性回归模型\textsuperscript{[13]}。该研究不仅做了显著性检验，还专门讨论了多重共线性和异方差问题，并进一步采用广义逐步回归法对模型进行改进，使改进后模型的拟合效果明显提升。其贡献在于为多因子选股提供了一套较规范的计量建模流程，但由于方法建立在二维回归框架上，因此难以同时处理股票、因子和时间三个维度的结构耦合。刘宇轩则进一步把宏观金融状态纳入多因子模型，在传统主成分分析法构建多因子选股模型的基础上，引入由 GARCH 模型合成的中国金融状况指数（FCI），构建基于金融周期的行业轮动多因子选股模型\textsuperscript{[14]}。该文不仅比较了金融周期指标与经济周期指标对选股模型的修正效果，还讨论了因子正负性会随市场状态变化而改变，并采用动态回归法判断当期因子正负性。其实证结果表明：引入宏观因素后，多因子模型收益表现得到改善，而金融周期指标相较于经济周期指标更能捕捉金融市场波动，从而取得更高的超额收益率。这说明传统多因子模型的一个关键不足，正是缺少对时间阶段与宏观状态的显式建模。

若把上述中文研究放在同一条问题链中观察，可以发现它们其实分别回答了几个彼此相关、但尚未被统一起来的问题。黄宏运和刘宇轩解决的是“给定因子集合后，如何利用统计回归、主成分分析和宏观修正提升选股效果”的问题\textsuperscript{[13][14]}；岳欣解决的是“当金融数据存在停牌缺失时，如何利用张量分解恢复多指标数据并开展后续分析”的问题\textsuperscript{[8]}；曾亚丽解决的是“当证券市场数据具有多市场、多指标和时间耦合特征时，如何通过张量特征提取提高预测精度”的问题\textsuperscript{[10]}；马少沛等解决的是“高阶张量降维时如何避免结构信息丢失并控制计算复杂度”的问题\textsuperscript{[9]}；李博则通过时态支持向量机和粒度变换，在股票操纵模式发现任务中说明了时间粒度变化本身就可能承载原始特征中难以直接发现的模式信息，其在未知数据集上的识别准确率可达到较高水平\textsuperscript{[18]}。这些研究分别触及了多因子选股、张量建模、时间结构识别和高维降维中的关键环节，但尚未将正式数据链、股票因子压缩、模式发现、排序有效性和样本边界稳健性评价统一进同一套研究框架。换言之，已有研究往往只覆盖“数据处理”“预测建模”或“因子选股”中的某一段链条，而本文希望补上的，正是这几段链条之间缺失的衔接关系。

综合来看，现有研究已经证明了张量分解在高维结构数据上的理论优势，也展示了其在金融时序、股票波动和相关结构建模中的应用潜力，但仍存在三个值得继续深入的问题。第一，专门针对“股票因子—多维结构保持降维—模式发现”的一体化研究仍不充分，尤其是面向中国 A 股样本的系统实证较少。第二，许多研究偏重于预测精度，较少同时讨论因子解释、股票联动和时间状态切换三类结果。第三，数据处理过程、实验结果与论文论证之间往往缺乏统一而清晰的证据对应关系\textsuperscript{[8][10][11][18]}。本文的研究正是围绕这些不足展开，尝试将张量分解方法、A 股因子数据与实证结果分析统一到同一研究框架中。

\section{研究内容与技术路线}

本文围绕“股票—因子—时间”三维张量构建研究对象，将研究内容划分为四个层面。第一，在数据层，构建可复用的 A 股正式数据链，完成 K 线、财务、报告和指数成分历史等基础数据的整理，并在此基础上形成因子面板。第二，在方法层，比较 CP 分解、Tucker 分解与 PCA 三种降维路径，重点分析张量方法在结构保持和潜在模式提取上的优势。第三，在结果层，围绕分解质量、排序有效性、时间状态变化和因子重要性展开实证分析。第四，在证据层，将数值结果、图形结果与文字论证统一到同一评价体系中，保证结论具有一致的证据来源。

对应的技术路线可以概括为“原始数据获取与整理—因子面板构建—三阶张量表示—张量分解与秩选择—潜在表示分析—实验评估与可视化输出”。其中，原始数据整理并不是附属步骤，而是后续张量建模能否成立的前提。由于股票市场数据天然存在停牌、成分股变化、财务披露时滞和不同数据源更新口径不一致等问题，本文在正式实验之前，先对 K 线前复权口径、财务/报告覆盖、指数成分历史和因子面板时间窗口进行统一，确保模型输入具有一致的时间跨度和样本边界。

在算法路线方面，本文先构造原始张量，再执行协议化预处理和训练/验证/测试切分，并分别在 CP、Tucker 和 PCA 三个模型上进行秩选择和留出评估。与仅比较重构误差的实验不同，本文同时关注三个问题：其一，哪种方法在重构质量上更优；其二，哪种方法在排序指标和滚动稳定性上表现更平衡；其三，哪种方法更有利于形成可解释的因子和时间模式。围绕这三个问题，正文给出表格、图形和定性讨论。

\section{论文结构安排}

除摘要、参考文献、附录等部分外，全文共分为六章。第一章为绪论，介绍研究背景、国内外研究现状、研究内容与技术路线。第二章阐述张量代数、张量分解与金融因子数据表示的理论基础。第三章给出数据来源、样本边界、预处理流程和张量构造方式。第四章详细说明基于张量分解的股票因子降维与模式发现方法。第五章基于正式实验结果，围绕分解质量、模式识别与排序有效性展开分析。第六章总结全文研究结论、创新点、不足与后续展望。

\chapter{理论基础与关键技术}

\section{张量代数基础}

在本文中，股票因子数据被表示为三阶张量$\mathcal{X}\in\mathbb{R}^{N\times F\times T}$，其中$N$表示股票数量，$F$表示因子数量，$T$表示观测时点数。与向量和矩阵相比，张量最大的特点在于能够同时保留多个模态的信息，因此更适合描述高维结构数据。若将张量简单展平成矩阵，则本属于不同模态的结构特征会被强行混合，造成解释上的模糊。基于此，本文直接在张量形式下进行建模。

为便于后文描述 CP 分解与 Tucker 分解，现对本文频繁使用的几类基础运算作出统一说明\textsuperscript{[4][6]}。首先，模态展开是指将张量某一模作为矩阵的行，其余模按固定顺序展开为列。对三阶张量$\mathcal{X}\in\mathbb{R}^{N\times F\times T}$而言，其三种展开矩阵分别记为$X_{(1)}\in\mathbb{R}^{N\times (FT)}$、$X_{(2)}\in\mathbb{R}^{F\times (NT)}$和$X_{(3)}\in\mathbb{R}^{T\times (NF)}$。该操作能够把高阶张量上的优化问题转化为矩阵形式，是后文交替最小二乘更新的计算基础。

其次，若向量$a\in\mathbb{R}^{N}$、$b\in\mathbb{R}^{F}$、$c\in\mathbb{R}^{T}$的外积记为$a\circ b\circ c$，则其生成一个秩一三阶张量，并满足
\begin{equation}
\left(a\circ b\circ c\right)_{ijt}=a_i b_j c_t.
\end{equation}
因此，CP 分解可以理解为若干个秩一张量的线性叠加。对于两个同型矩阵$M,N\in\mathbb{R}^{I\times J}$，其 Hadamard 积记为$M\ast N$，定义为
\begin{equation}
\left(M\ast N\right)_{ij}=m_{ij}n_{ij}.
\end{equation}
若$A=[a_1,\ldots,a_R]\in\mathbb{R}^{I\times R}$、$B=[b_1,\ldots,b_R]\in\mathbb{R}^{J\times R}$，则 Khatri--Rao 积记为
\begin{equation}
A\odot B=\left[a_1\otimes b_1,\;a_2\otimes b_2,\;\ldots,\;a_R\otimes b_R\right]\in\mathbb{R}^{(IJ)\times R},
\end{equation}
其中$\otimes$表示 Kronecker 积。Khatri--Rao 积与 Hadamard 积将在后文的 CP 更新公式中直接出现。

最后，张量与矩阵的$n$模积是 Tucker 分解的核心运算。设$\mathcal{X}\in\mathbb{R}^{I_1\times\cdots\times I_n\times\cdots\times I_N}$，且$U\in\mathbb{R}^{J\times I_n}$，则$\mathcal{X}\times_n U$表示$\mathcal{X}$沿第$n$模与$U$相乘，其结果为
\begin{equation}
\left(\mathcal{X}\times_n U\right)_{i_1\cdots i_{n-1}ji_{n+1}\cdots i_N}
=\sum_{i_n=1}^{I_n}x_{i_1\cdots i_n\cdots i_N}u_{ji_n}.
\end{equation}
由此可见，模态展开、外积、Khatri--Rao 积、Hadamard 积以及$n$模积分别对应了本文后续 CP 分解和 Tucker 分解中的关键符号与计算步骤。

在股票因子应用中，张量代数的价值不仅体现在数学表示上，也体现在金融解释上。股票模态对应横截面结构，因子模态对应信息维度，时间模态对应市场演化。若某一潜在成分在股票模态上的载荷主要集中于银行、券商和保险股，而在时间模态上又集中出现在政策宽松阶段，则该潜在成分就可能对应“金融风格因子在特定市场环境下的整体共振”。因此，张量的每个模态都具有明确的经济含义，而不是单纯的索引维度。

\section{张量分解模型}

张量分解的基本思想，是在尽可能保留原始多模态结构的前提下，用较少的潜在成分对高维张量进行近似表示。对本文的股票因子张量而言，这一过程既对应因子压缩，也对应股票结构与时间状态的提取。综合经典理论与本文研究目标，后文主要使用 CP 分解与 Tucker 分解两类模型\textsuperscript{[1][2][3][4]}。

\subsection{CP 分解}

CP 分解将原始张量表示为若干个秩一张量之和，其数学形式为
\begin{equation}
\mathcal{X}\approx \sum_{r=1}^{R}\lambda_r a_r\circ b_r\circ c_r.
\end{equation}
其中，$R$为分解秩，$\lambda_r$为第$r$个成分的权重，$a_r$、$b_r$、$c_r$分别表示第$r$个成分在股票、因子和时间三个模态上的载荷向量。若将这些向量按列组合，可得因子矩阵$A=[a_1,\ldots,a_R]$、$B=[b_1,\ldots,b_R]$和$C=[c_1,\ldots,c_R]$。从经济含义上看，矩阵$A$描述股票在潜在空间中的位置，矩阵$B$反映原始因子对潜在成分的贡献强弱，矩阵$C$则揭示潜在结构随时间的变化轨迹。

CP 分解对应的目标可以写为
\begin{equation}
\min_{A,B,C,\lambda}\left\|\mathcal{X}-\sum_{r=1}^{R}\lambda_r a_r\circ b_r\circ c_r\right\|_F^2.
\end{equation}
该模型的优点在于结构简洁、参数量相对可控，并且当唯一性条件成立时，分解结果具有较好的可解释性\textsuperscript{[4][5]}。在算法实现中，本文采用交替最小二乘思想逐模更新因子矩阵。以固定$B$和$C$时对$A$的更新为例，可写为
\begin{equation}
A\leftarrow X_{(1)}(C\odot B)\left[(C^{\top}C)\ast(B^{\top}B)\right]^{-1},
\end{equation}
其中$\odot$表示 Khatri--Rao 积，$\ast$表示 Hadamard 积。由此可以看出，CP 分解实质上是通过模态展开与矩阵优化，逐步逼近原始张量的低秩结构。

\subsection{Tucker 分解}

与 CP 分解不同，Tucker 分解将原始张量表示为核张量与多个因子矩阵的$n$模乘积，即
\begin{equation}
\mathcal{X}\approx \mathcal{G}\times_1 A\times_2 B\times_3 C.
\end{equation}
其中，$\mathcal{G}\in\mathbb{R}^{R_1\times R_2\times R_3}$为核张量，$A\in\mathbb{R}^{N\times R_1}$、$B\in\mathbb{R}^{F\times R_2}$、$C\in\mathbb{R}^{T\times R_3}$分别对应三个模态上的低维表示矩阵。核张量$\mathcal{G}$不再要求对角结构，因此能够刻画不同潜在成分之间更复杂的高阶交互关系。

Tucker 分解对应的优化目标可写为
\begin{equation}
\min_{\mathcal{G},A,B,C}\left\|\mathcal{X}-\mathcal{G}\times_1 A\times_2 B\times_3 C\right\|_F^2.
\end{equation}
该模型最重要的特点，是允许不同模态采用不同秩$(R_1,R_2,R_3)$。例如，在股票模态保留较高维度以刻画横截面差异，而在因子模态和时间模态保留较低维度以提炼主导结构。这种多线性秩约束使 Tucker 分解在面对非对称、强耦合的金融数据时更具灵活性。求解时，通常先利用 HOSVD 或截断奇异值分解获得初始因子矩阵，再通过交替迭代更新核张量与各模态因子矩阵。

\subsection{两类模型的比较与本文选择}

从结构上看，CP 分解可视为 Tucker 分解在核张量退化为超对角结构时的特例，因此两者都属于低秩张量表示框架，但强调的建模能力并不相同。CP 分解更强调用少量共享成分解释整体结构，适合观察是否存在较稳定、较统一的潜在模式；Tucker 分解则通过核张量保留成分之间的交互关系，更适合描述股票、因子与时间三者之间并不对称的复杂耦合。

就本文研究目标而言，CP 分解能够提供一个结构清晰、便于比较的低秩基线；Tucker 分解则更适合承担主体分析任务，因为股票因子数据通常同时包含行业分化、风格轮动和阶段切换等多层次变化，仅靠秩一张量叠加往往难以完整刻画全部结构。因此，后文将同时报告 CP 与 Tucker 的结果，并以 PCA 作为二维降维对照，从重构能力、排序有效性和结构解释性三个维度比较不同方法的适用边界。

\section{金融数据的张量表示方法}

在股票因子问题中，最自然的张量表示方式是以股票为第一模、因子为第二模、时间为第三模，将每个元素$x_{ijt}$定义为第$i$只股票在第$t$个交易日上的第$j$个因子取值，即
\begin{equation}
\mathcal{X}=(x_{ijt})\in\mathbb{R}^{N\times F\times T}.
\end{equation}
这一表示方式具有三方面优势。第一，它保留了股票横截面之间的相对关系，可直接观察行业、风格和规模层面的结构接近性。第二，它保留了因子之间的共振关系，适合识别冗余因子和具有替代关系的因子簇。第三，它保留了时间方向上的连续性和阶段性，能够描述因子结构在牛熊转换、流动性冲击和政策扰动中的演化。

若进一步考虑正式实验中的样本池边界，则同样可以在统一结构下表示 HS300、SZ50 与 ZZ500 三个样本池对应的因子张量，只需在样本选择阶段通过指数成分历史对股票集合进行限定。由于本文目标并非构建一个全市场实时交易系统，而是完成结构保持降维与模式发现，因此更关注样本口径一致、时间窗口完整和特征构造稳定，而不是在不同样本池之间混合建模。

\section{模型评价指标}

为了评价降维效果，本文使用重构误差、解释方差、排序相关、信息比率和滚动稳定性等指标。首先，重构误差衡量分解后张量与原始张量之间的偏离程度，可写为
\begin{equation}
\mathrm{MSE}=\frac{1}{NFT}\|\mathcal{X}-\hat{\mathcal{X}}\|_F^2.
\end{equation}
与之对应，解释方差反映模型对原始张量波动信息的保留程度，数值越大通常表示降维后保留的有效结构越多。

其次，考虑到本文不仅关注重构能力，也关注低维表示的排序价值，因此还使用 Rank IC 和 IR 指标来观察降维后的潜在得分与下一期收益之间的相关关系。其中，Rank IC 侧重于秩相关，适合金融场景下强调排序而非绝对值预测的任务。若某模型在解释方差上较高但 Rank IC 很低，则意味着它更擅长数值重构而不一定更适合选股排序。

最后，本文引入滚动稳定性指标衡量模型在不同时间窗口下输出潜在结构的一致性。对因子降维而言，若模型在不同窗口中频繁产生完全不同的潜在结构，则即使单次拟合表现良好，也很难形成可复用的研究结论。因此，重构误差、解释方差、Rank IC 与滚动稳定性共同构成本文的评价体系，其中前两者侧重数值拟合，后两者侧重金融解释与研究稳健性。

\chapter{数据预处理与张量表示}

\section{数据来源与样本说明}

本文使用的数据主要来源于自建的 A 股研究数据库。原始信息以 Baostock 数据为主，并进一步整理为日频 K 线、财务数据、业绩报告、指数成分历史、分红复权信息和衍生因子面板等多类表格。本文选取的正式样本池分别为 HS300、SZ50 和 ZZ500。研究设计上，希望观察的时间跨度覆盖 2015 年 1 月至 2026 年 4 月，以便同时考察多个市场阶段；就本文当前纳入正文比较的结果而言，核心短窗口样本覆盖 2026 年 3 月 2 日至 2026 年 3 月 30 日。之所以仍以这三个样本池为主，是因为它们具有较清晰的风格边界、较好的流动性基础和较稳定的成分定义，能够更清楚地观察因子结构与样本池风格之间的对应关系。

从数据覆盖角度看，当前研究数据已经完成前复权 K 线、财务表、报告表和指数成分历史的整理，但原始层数据的可用跨度与当前正文重点讨论的短窗口实验之间仍然存在差异。也就是说，数据基础具备继续向长期窗口扩展的条件，而本文主要展示的比较结果仍集中于 2026 年 3 月的短窗口样本。这也是本文在局限性分析中专门强调“实验窗口口径仍需进一步统一”的原因。

需要强调的是，“数据已经抓取完成”与“数据已经适合直接进入模型训练”并不是同一概念。当前训练张量仍以主线因子面板为核心，但扩展方案已经在第一阶段市场代理变量、财务 PIT 指标、业绩快报与业绩预告事件特征之外，进一步纳入外部利率、月度宏观指标、分红信息、重大事项公告与公告标题文本。为了避免在训练阶段引入披露滞后和未来信息泄露，本文对这些扩展特征的披露时点与可用时点进行了单独梳理，并将其划分为市场代理、外部宏观、财务点时特征、分红公司行为和公告事件字典五类。各类特征均通过 \texttt{pub\_date} 与 \texttt{available\_date} 的映射规则进入训练样本，这一安排既解释了当前输入边界的形成方式，也为后续继续扩展更细粒度源表提供了依据。

在样本组织上，本文不直接把所有股票合并到同一张量中统一训练，而是分别对 HS300、SZ50 和 ZZ500 构建因子张量。这种做法的动机在于，不同指数样本池本身具有明显不同的规模特征、行业构成和风格偏向。若直接混合，潜在表示可能首先学习到“样本池差异”，而非股票因子内部的结构关系。分样本池建模则更有利于在同质性更高的截面内观察张量分解的表现。

\section{股票池与时间窗口设定}

指数样本池的边界并非静态固定，因此本文采用“指数成分历史”而非单一时点快照定义股票池。对于任一交易日，只将当日属于目标指数成分集合的股票纳入样本。这样做有两个好处：第一，避免了使用当前成分股去回溯历史样本而产生的未来信息泄露；第二，能够更真实地反映指数样本池在长期窗口中的构成变化。

从研究设计角度看，更理想的正式窗口应覆盖 2015 年 1 月至 2026 年 4 月，以便同时观察多轮市场阶段变化；但本文正文的核心基线比较仍以 2026 年 3 月的短窗口结果为主。因此，后文中首先讨论的实证结论应理解为短窗口口径下的结果，而不是对长期全窗口的直接替代。

\section{数据预处理}

数据预处理的目标，是在不破坏“股票—因子—时间”原始结构的前提下，将正式因子面板整理为可直接进入张量分解模型的稳定输入。相较于仅做一次静态清洗，本文更强调“先对齐样本边界，再做缺失填补、异常值收缩和标准化，最后按时间顺序切分训练、验证与测试样本”的协议化处理流程。表\ref{tab:preprocess-panel-summary}给出了本文三类正式样本池在进入模型前的面板规模，可以看到三个样本池均已整理为无缺失的完整短窗口因子面板。

\begin{table}[htbp]
\centering
\caption{正式样本池因子面板规模与完整性}
\label{tab:preprocess-panel-summary}
\small
\begin{tabular}{lccccc}
\toprule
样本池 & 交易日数 & 股票数 & 因子数 & 因子观测值总数 & 缺失率 \\
\midrule
HS300 & 21 & 72 & 7 & 10\,584 & 0.00\% \\
SZ50  & 21 & 50 & 7 & 7\,350  & 0.00\% \\
ZZ500 & 21 & 500 & 7 & 73\,500 & 0.00\% \\
\bottomrule
\end{tabular}
\end{table}

从处理顺序看，本文采用的预处理流程如图\ref{fig:preprocess-flow}所示。该流程先保证样本口径与交易日对齐，再逐步完成缺失值填补、异常值收缩和标准化，最后输出三阶张量。这样做的目的，是避免不同预处理步骤相互穿插而造成口径混乱。

\begin{figure}[htbp]
\centering
\begin{tikzpicture}[x=1cm,y=1cm]
  \node[draw, rounded corners, align=center, minimum width=3.1cm, minimum height=1cm] (a) at (0,0) {原始因子面板\\统一股票池与交易日};
  \node[draw, rounded corners, align=center, minimum width=3.1cm, minimum height=1cm] (b) at (3.8,0) {缺失值处理\\前向填充/后向填充/行业中位数};
  \node[draw, rounded corners, align=center, minimum width=3.1cm, minimum height=1cm] (c) at (7.6,0) {异常值收缩\\截尾或 MAD 稳健收缩};
  \node[draw, rounded corners, align=center, minimum width=3.1cm, minimum height=1cm] (d) at (1.9,-2.2) {训练集标准化\\验证集与测试集复用参数};
  \node[draw, rounded corners, align=center, minimum width=3.1cm, minimum height=1cm] (e) at (5.7,-2.2) {按时间顺序切分\\训练 / 验证 / 测试并构造张量};
  \draw[->, thick] (a) -- (b);
  \draw[->, thick] (b) -- (c);
  \draw[->, thick] (c) |- (d);
  \draw[->, thick] (d) -- (e);
\end{tikzpicture}
\caption{正式实验中的因子预处理流程}
\label{fig:preprocess-flow}
\end{figure}

首先，在缺失值处理阶段，股票因子缺失往往来自停牌、上市时间不足、财务数据尚未披露或数据源覆盖限制。若直接删除含缺失值样本，会造成时间窗口破碎和不同样本池之间的不平衡，因此本文采用分层填补规则。设原始因子值为$x_{ijt}$，则填补后的值$\tilde{x}_{ijt}$可写为
\begin{equation}
\tilde{x}_{ijt}=
\begin{cases}
x_{ijt}, & x_{ijt}\text{ 已观测};\\
x_{ij,t^-}, & x_{ijt}\text{ 缺失且存在最近一期历史值};\\
x_{ij,t^+}, & x_{ijt}\text{ 缺失且不存在历史值但存在后续值};\\
\operatorname{med}_{g(i),j,t}, & \text{以上条件均不满足},
\end{cases}
\end{equation}
其中$g(i)$表示股票$i$所属行业，$\operatorname{med}_{g(i),j,t}$表示该行业在时点$t$上因子$j$的截面中位数。该规则兼顾了时间连续性与行业可比性。

其次，在异常值处理阶段，估值类、换手类和未来收益类因子容易受极端行情或低基数效应影响。为避免张量分解被少量极端值主导，本文以训练集统计量为基准，对因子进行稳健收缩。设$\operatorname{med}_j$和$\operatorname{MAD}_j$分别表示因子$j$在训练集中的中位数与中位数绝对偏差，则收缩后的因子值$\hat{x}_{ijt}$写为
\begin{equation}
\hat{x}_{ijt}
=\min\left\{\max\left(\tilde{x}_{ijt},\operatorname{med}_j-3\operatorname{MAD}_j\right),\operatorname{med}_j+3\operatorname{MAD}_j\right\}.
\end{equation}
这一处理能够在保留排序信息的同时抑制异常跳点对 Frobenius 范数的放大效应。

再次，在标准化阶段，本文仅使用训练集估计位置与尺度参数，再将同一组参数作用于验证集和测试集，以避免未来信息泄露。设$\mu_j^{\mathrm{train}}$和$\sigma_j^{\mathrm{train}}$分别表示因子$j$在训练集中的均值和标准差，则标准化结果$z_{ijt}$定义为
\begin{equation}
z_{ijt}=\frac{\hat{x}_{ijt}-\mu_j^{\mathrm{train}}}{\sigma_j^{\mathrm{train}}+\varepsilon},
\end{equation}
其中$\varepsilon$为防止分母为零而引入的极小常数。经过该步骤后，不同量纲的因子被映射到可比较的尺度上，从而避免某一类数值较大的因子在分解过程中占据过高权重。

最后，在样本切分与边界控制方面，本文不采用“全样本一次性拟合”的做法，而是显式划分训练、验证和测试区间：训练集负责拟合预处理状态和模型参数，验证集用于选择秩，测试集仅用于最终 held-out 评估。对于扩展特征，只有披露时点与可交易日映射规则已经明确的数据才进入训练输入，其余时点定义尚不稳定的补充信息暂不并入主线张量。通过上述安排，本文尽可能保证了实验口径、数据边界与模型评价之间的一致性。

\section{张量表示}

在完成预处理之后，本文将每个样本池的数据表示为三阶张量$\mathcal{X}\in\mathbb{R}^{N\times F\times T}$。其中，$N$是该样本池在给定交易日有效成分股集合下的股票数量，$F$是因子面板中的因子数，$T$是时间维度上的交易日数。由于指数成分历史会导致股票集合在不同时间上变化，工程实现中采用了统一编码与缺失补位的方式，以保证张量在模态维度上保持固定形状。

这种张量表示与传统矩阵表示的最大区别在于，时间维度不再只是“样本编号的一部分”，而是与股票和因子维度共同组成潜在结构的核心模态。由此得到的模型结果不仅包含重构误差等数值信息，也会输出三个模态上的潜在表示矩阵，从而为后续的股票相似性、因子重要性和时间状态变化分析提供直接依据。

\chapter{基于张量分解的股票因子降维与模式发现方法}

\section{问题定义}

本文研究的问题可以形式化为：给定股票—因子—时间三阶张量$\mathcal{X}$，如何在保持其多维结构信息的前提下，构造一个低维表示$\hat{\mathcal{X}}$，使得其既能较好重构原始张量，又能输出对股票、因子与时间具有解释意义的潜在表示。相比一般的“仅降维”问题，本文更强调“降维+模式发现”的双重目标。

具体而言，本文希望完成以下三个子任务。第一，在因子模态上实现有效压缩，从而降低高维因子集合中的冗余程度。第二，在股票模态上识别潜在的相似结构，用于解释哪些股票在潜在风格暴露上更接近。第三，在时间模态上识别市场阶段变化，观察潜在结构在长期样本中的迁移与突变。围绕这三个目标，单纯使用二维 PCA 虽然可以提供重构基线，但不具备自然的多模态解释能力，因此需要引入张量分解方法。

\section{张量分解模型构建}

本文同时构建 CP、Tucker 与 PCA 三类模型。CP 模型作为最简洁的张量分解基线，用于观察在低秩约束较强时，三阶张量能否被少量公共成分有效表示。Tucker 模型作为论文主方法，用于在不同模态设置不同秩，捕捉更复杂的模态间交互结构。PCA 模型虽然不在张量形式下建模，但由于其在传统降维问题中具有较强数值表现，因此作为二维压缩基线保留。

对 CP 分解而言，模型秩$R$越大，理论上可表示的结构越丰富，但解释性和压缩率会下降；秩过小则可能无法表示足够的结构信息。对 Tucker 分解而言，需要同时设置股票模态秩、因子模态秩和时间模态秩。本文在实验中并不依赖人工经验一次性固定秩，而是先给出候选 rank 集合，再在验证集上比较不同 rank 对应的 held-out MSE，选出表现最优的参数组合。这种做法能在一定程度上避免“为了论文结论而人为挑选 rank”的主观性问题。

在因子筛选与模式解释上，本文主要依赖因子模态矩阵$B$及其载荷大小。若某个潜在成分在$B$中对质量因子、估值因子和动量因子分别给出不同权重，则说明不同金融含义的因子被组合进不同潜在方向中。进一步地，可定义因子$j$的综合重要性为
\begin{equation}
w_j=\frac{\|b_{j:}\|_2}{\sum_{k=1}^{F}\|b_{k:}\|_2},
\end{equation}
并按$w_j$排序识别核心因子。这一做法虽然不直接等价于经济学中的因子定价显著性检验，但对于“哪些因子在低维潜在结构中承担了主要信息功能”这一问题具有直观解释力。

\section{模型求解与实现}

本文的实验流程可以概括为五步。第一，读取样本池对应的因子面板和指数历史，构造原始三阶张量。第二，按照显式训练/验证/测试协议执行预处理和切分。第三，对每一类模型分别在候选秩集合上进行训练和验证，选择最优参数。第四，在训练集与验证集联合数据上重估最终模型，并只在测试集上进行留出评估。第五，输出模型评价指标、因子重要性、时间状态、候选股票排序以及相应图形结果。

在算法层面，CP 和 Tucker 模型均依赖迭代更新。为了减轻局部最优对结果的影响，本文使用固定随机种子和稳定的预处理协议，并保留完整的秩选择与评价记录。PCA 则作为二维基线，以相同测试集口径进行评估，确保三类方法比较公平。需要指出的是，本文关注的是研究型实证框架，而不是高频交易系统，因此方法实现更强调可复现性、可追溯性和结果比较的一致性，而非极限运行速度。

\chapter{实验结果与分析}

\section{实验设置}

正式实验样本池包括 HS300、SZ50 与 ZZ500。本文正文首先比较 2026 年 3 月 2 日至 2026 年 3 月 30 日窗口下的正式结果。每个样本池均基于正式因子面板构建张量，并在统一的预处理、切分和评价协议下比较 CP、Tucker 与 PCA 三种方法。其中，CP 和 Tucker 负责提供张量结构建模结果，PCA 作为二维压缩基线用于比较其数值拟合能力。

在组合评价口径上，本文统一采用 A 股 T+1 交易假设：以 $t$ 日可观测因子形成排序信号，于 $t+1$ 日建仓，并按照 5 个交易日持有期计算远期收益；交易成本按单边万分之 1 与换手率线性对应扣减。这样的设定既避免了当日信号当日成交的前视偏差，也使组合层结果更接近 A 股市场的实际执行约束。

从输出结果看，实验不仅生成数值指标，还形成了解释方差对比图、Rank IC 对比图、模型指标总览图、时间状态变化图、因子重要性热力图，以及组合累计净值对比图和回撤曲线对比图。这样的图文并行呈现方式有助于把论文从“结果表格堆砌”转向“结构模式分析”，也是本文更强调的一点。

\section{分解质量分析}

首先从数值重构能力出发，三组样本池上的结果具有较清晰的一致性。HS300 样本中，CP、Tucker、PCA 的 MSE 分别为 0.9140、0.3406 和接近 0，对应解释方差分别为 0.0017、0.6280 和 1.0000；SZ50 样本中，三者的 MSE 分别为 0.9898、0.3160 和接近 0，对应解释方差分别为 $-0.0527$、0.6639 和 1.0000；ZZ500 样本中，三者的 MSE 分别为 1.0091、0.4832 和接近 0，对应解释方差分别为 $-0.0541$、0.4953 和 1.0000。可以看到，在当前短窗口结果下，PCA 在三组样本池上都表现出最强的重构能力，而 Tucker 在两个张量方法中稳定优于 CP。

这一结果说明，若只以数值重构为目标，二维压缩方法仍具有明显优势。原因在于，PCA 在矩阵展开后直接针对方差最大方向建模，天然更偏向“数值逼近最优”；而张量方法在保持多模态结构约束的同时，必须在不同模态间平衡表示能力，因此纯重构性能未必领先。然而，这并不意味着张量方法失去研究价值。恰恰相反，若研究目标是结构解释与模式发现，那么在重构精度略有让步的情况下换取更清晰的多模态潜在表示，是完全合理的策略。

\begin{figure}[!t]
\centering
\begin{tikzpicture}
\begin{groupplot}[
  group style={group size=3 by 1, horizontal sep=1.3cm},
  width=0.28\textwidth,
  height=6.0cm,
  ybar,
  ymin=-0.1,
  ymax=1.08,
  symbolic x coords={CP,Tucker,PCA},
  xtick=data,
  enlarge y limits={upper=0.06,lower=0.04},
  extra y ticks={0},
  extra y tick labels={},
  extra y tick style={grid=major, major grid style={black!65, line width=0.35pt}},
  point meta=y,
  nodes near coords={\pgfmathprintnumber[fixed,precision=2]{\pgfplotspointmeta}},
  nodes near coords,
  nodes near coords align={vertical},
  every node near coord/.append style={font=\tiny},
  tick label style={font=\small},
  ylabel style={font=\small},
  title style={font=\small},
  legend style={font=\small, draw=none}
]
\nextgroupplot[title={HS300},ylabel={解释方差}]
\addplot[fill=blue!55] coordinates {(CP,0.0017) (Tucker,0.6280) (PCA,1.0000)};
\nextgroupplot[title={SZ50}]
\addplot[fill=orange!70] coordinates {(CP,-0.0527) (Tucker,0.6639) (PCA,1.0000)};
\nextgroupplot[title={ZZ500}]
\addplot[fill=teal!65] coordinates {(CP,-0.0541) (Tucker,0.4953) (PCA,1.0000)};
\end{groupplot}
\end{tikzpicture}
\caption{三个正式样本池上的模型解释方差对比图}
\label{fig:variance}
\end{figure}

图~\ref{fig:variance} 更直观地展示了解释方差差异。PCA 在三个样本池上均显著领先；Tucker 稳定保持较高的正解释方差，而 CP 只有在 HS300 上勉强保持接近于零的解释方差，在 SZ50 和 ZZ500 上仍为负值。这表明，在当前短窗口口径内，低秩 CP 仍不足以稳定表示股票因子结构，而 Tucker 在更灵活的多线性秩约束下能够更稳定地提取结构信息。

\section{模式发现结果分析}

除了数值重构能力外，本文更关注潜在表示所揭示的模式结构。根据因子重要性输出，价值因子、质量因子和波动率因子在多个样本池和多个模型下都反复出现于前列，说明它们构成了当前正式因子面板中的核心信息轴。以当前 HS300 的 Tucker 结果为例，波动率因子、价值因子和质量因子的平均重要性依次居前，而动量因子在当前短窗口内的重要性接近于零，说明在这一窗口下市场结构更偏向估值、质量与波动信息，而非长周期趋势信息。

在股票模态层面，潜在表示可用于构造股票相似关系。虽然本文不把股票相似矩阵单独作为最终评价指标，但从实验日志和相似对输出可以观察到，同一行业或相近风格股票更容易在潜在空间中表现出接近的载荷结构。这与量化投资中“因子暴露决定风格分组”的直觉是一致的。相较于简单相关系数，张量潜在表示构造的相似关系同时考虑了时间和因子两个方向，因此更容易反映结构性接近，而非仅仅是短期价格同步。

进一步基于全 A 样本的 Tucker 结果，可以得到股票潜在结构图与聚类—行业交叉图。前者将股票按潜在聚类标签展开，并用行业颜色展示聚类内部的行业构成；后者则把聚类标签与行业分布做交叉统计。这两幅图共同说明，股票潜在空间并不是简单复制行业分类，而是在行业信息之外保留了由因子暴露、时间状态和排序得分共同形成的结构相似性。

\begin{figure}[!t]
\centering
\begin{tikzpicture}
\begin{groupplot}[
  group style={group size=2 by 1, horizontal sep=2.2cm},
  width=0.36\textwidth,
  height=7.0cm,
  xbar,
  enlarge y limits=0.18,
  yticklabel style={font=\scriptsize, text width=2.0cm, align=right},
  tick label style={font=\small},
  label style={font=\small},
  title style={font=\small},
  point meta=x,
  nodes near coords={\pgfmathprintnumber[fixed,precision=2]{\pgfplotspointmeta}},
  nodes near coords,
  every node near coord/.append style={font=\tiny},
]
\nextgroupplot[
  title={Tucker跨边界Rank IC},
  xmin=-0.42,
  xmax=0.28,
  symbolic y coords={大市值,小市值,电子C39,专用设备C35,中市值,中证500,全市场,沪深300,上证50,医药C27},
  ytick=data,
  xlabel={Rank IC}
]
\addplot[fill=blue!55] coordinates {
(-0.3780,大市值)
(-0.0997,小市值)
(-0.1196,电子C39)
(-0.0496,专用设备C35)
(0.1031,中市值)
(0.0209,中证500)
(0.0279,全市场)
(-0.0468,沪深300)
(0.0468,上证50)
(0.1476,医药C27)
};
\nextgroupplot[
  title={Tucker跨边界累计收益},
  xmin=-12,
  xmax=32,
  symbolic y coords={大市值,小市值,电子C39,中市值,全市场,中证500,上证50,专用设备C35,沪深300,医药C27},
  ytick=data,
  xlabel={累计收益(\%)}
]
\addplot[fill=orange!70] coordinates {
(-7.26,小市值)
(-3.65,电子C39)
(-0.28,大市值)
(-2.11,全市场)
(-1.63,中证500)
(-0.43,中市值)
(0.04,上证50)
(0.17,专用设备C35)
(1.29,沪深300)
(16.46,医药C27)
};
\end{groupplot}
\end{tikzpicture}
\caption{短窗口样本边界扩展下 Tucker 的跨边界排序与收益分化}
\label{fig:tucker_boundary_short}
\end{figure}

时间模态则提供了对市场状态切换的观察窗口。HS300 的 Tucker 结果显示，最明显的相邻日期结构跳变集中在 2026 年 3 月 25 日至 2026 年 3 月 30 日这一短区间，说明即使在较短的正式输出窗口内，潜在结构也并非静态稳定，而是在局部阶段出现明显重排。换言之，张量分解不仅用于“压缩因子”，还能够识别“市场何时发生了潜在结构变化”。

\begin{figure}[!t]
\centering
\begin{tikzpicture}
\begin{groupplot}[
  group style={group size=3 by 1, horizontal sep=1.3cm},
  width=0.28\textwidth,
  height=6.0cm,
  ybar,
  ymin=-0.10,
  ymax=0.18,
  symbolic x coords={CP,Tucker,PCA},
  xtick=data,
  enlarge y limits={upper=0.10,lower=0.08},
  enlarge x limits=0.18,
  point meta=y,
  nodes near coords={\pgfmathprintnumber[fixed,precision=2]{\pgfplotspointmeta}},
  nodes near coords,
  nodes near coords align={vertical},
  every node near coord/.append style={font=\tiny},
  tick label style={font=\small},
  ylabel style={font=\small},
  title style={font=\small}
]
\nextgroupplot[title={HS300},ylabel={Rank IC}]
\addplot[fill=blue!55] coordinates {(CP,-0.0278) (Tucker,-0.0468) (PCA,-0.0530)};
\nextgroupplot[title={SZ50}]
\addplot[fill=orange!70] coordinates {(CP,0.0198) (Tucker,0.0468) (PCA,-0.0224)};
\nextgroupplot[title={ZZ500}]
\addplot[fill=teal!65] coordinates {(CP,-0.1082) (Tucker,0.0209) (PCA,-0.0442)};
\end{groupplot}
\end{tikzpicture}
\caption{三个正式样本池上的模型 Rank IC 对比图}
\label{fig:rankic}
\end{figure}

图~\ref{fig:rankic} 展示了三个正式样本池在 Rank IC 指标上的差异结构。可以看到，在严格 T+1 口径下，HS300 的三类模型 Rank IC 均转为负值，但幅度仍集中在较低区间；SZ50 与 ZZ500 中则分别只有 Tucker 保持了 0.0468 和 0.0209 的正向排序相关。这意味着仅有较强重构能力并不足以保证排序有效性，保持三维结构的潜在表示在部分样本池中仍具有独立价值。

\section{选股有效性分析}

为了从金融意义上评价降维结果，本文进一步考察 Rank IC 指标。图~\ref{fig:rankic} 显示，在 HS300 样本中，CP、Tucker 与 PCA 的 Rank IC 均值分别为 $-0.0278$、$-0.0468$ 与 $-0.0530$，说明大盘蓝筹样本在严格 T+1 口径下的短窗口排序难度更高；在 SZ50 样本中，Tucker 为 0.0468、CP 为 0.0198，而 PCA 回落到 $-0.0224$；在 ZZ500 样本中，也只有 Tucker 维持了 0.0209 的正向排序相关。可见，若从“跨样本池保持正向排序相关”这一视角出发，Tucker 在当前正式窗口内仍然表现最为均衡。

这一结果说明，数值拟合最强的方法不一定也是排序最有效的方法。PCA 在解释方差上显著领先，但在 ZZ500 上的 Rank IC 仍为负值，意味着其压缩方向更偏向静态方差解释，而不一定适合直接服务于股票横截面排序。相反，Tucker 通过保持三维结构，在因子压缩的同时保留了更多与时间演化和股票截面相关的模式，因此在三个样本池上都没有出现 Rank IC 由正转负的情况。

从选股研究角度看，这一结论具有现实含义。实际量化研究并不一定要求低维表示精确重构每一个原始因子值，而是更关注压缩后的潜在表示能否支持稳健的股票排序与风格解释。若一个模型在重构误差上最优、但排序效果不稳定，则它更接近数值压缩工具；若一个模型在排序指标和滚动稳定性上都更均衡，则说明其潜在表示更有利于后续金融解释与截面决策。当前结果表明，在短窗口正式比较中，Tucker 在这一点上更具优势。

进一步把排序结果转化为组合层指标后，可以观察到“排序相关性”和“组合表现”并不完全同步。在基准方案结果中，HS300 上三类模型组合都取得了正累计收益，其中 Tucker 达到 1.29\%，高于 PCA 的 0.92\% 和 CP 的 0.53\%；SZ50 上只有 Tucker 保持微弱正收益 0.04\%，而 CP 与 PCA 分别回落到 $-2.27\%$ 和 $-0.56\%$；ZZ500 上 Tucker 与 PCA 的累计收益分别为 $-1.63\%$ 和 $-1.82\%$，虽然都未转正，但仍显著好于 CP 的 $-7.76\%$。换言之，单纯较高的 Rank IC 并不自动等价于更优的组合净值，这也是将组合结果与 Rank IC 联动分析的必要性所在。

在扩展方案结果中，这种差异性更加明显。HS300 上扩展特征使三类模型的组合表现同步恶化，CP、Tucker 和 PCA 的累计收益分别降至 $-1.88\%$、$-3.12\%$ 和 $-2.56\%$，对应 Rank IC 也全部转为负值；SZ50 上扩展输入同样未改善排序表现，CP、Tucker 与 PCA 的累计收益分别为 $-3.07\%$、$-4.12\%$ 与 $-3.49\%$；ZZ500 上则出现分化，CP 的累计收益转为 1.18\%，同时 Rank IC 升至 0.0342，而 Tucker 与 PCA 仍分别回落到 $-3.39\%$ 和 $-3.47\%$。这表明，扩展特征不仅会改变模型的重构与排序结果，也会改变它们在组合层的收益—风险权衡，因此在评价扩展输入时，必须同时观察净值曲线、回撤、换手率和行业/风格暴露，而不能只保留单一排序指标。

在进一步补充分位数组、多空组合、成本后净值和基准超额收益之后，可以更细致地判断“排序有效性是否真正能够转化为可配置收益”。例如，基准方案下的 HS300 样本中，三类模型的 \texttt{top-bottom} 多空净值全部为正，其中 PCA、Tucker 与 CP 的多空累计净值分别达到 1.85\%、1.71\% 和 1.68\%；但一旦计入交易成本，三类模型的成本后累计收益分别回落到 0.91\%、1.28\% 和 0.51\%，说明毛收益与换手约束并不总是同向变化。又如基准方案下的 SZ50 样本中，Tucker 的多空组合累计净值达到 7.50\%，成本后净值仍基本持平，并保留了 0.04\% 的微弱基准超额，说明其分位排序结构虽弱但并未被基准完全吞没。扩展方案下的 ZZ500 样本中，Tucker 与 PCA 的成本后净值分别回落到 $-3.40\%$ 和 $-3.48\%$，相对基准的超额净值也分别为 $-5.40\%$ 和 $-5.48\%$，说明更宽输入边界虽然改变了结构解释，但在当前短窗口中并未把中盘样本稳定转化为可验证的正向配置价值。可见，加入分位数组、多空、成本和超额收益之后，组合层证据已经比单纯的前N净值曲线更接近真实投资决策场景。

为了把组合层结果与样本边界、行业暴露和风格暴露放到同一视角下比较，本文进一步构造了统一的跨边界汇总表。该表同时列出各样本池的 Rank IC 最优模型、累计收益最优模型、Sharpe 最优模型，以及每个模型的前N组合收益、分组价差、多空净值、成本后净值、平均交易成本、超额净值、年化波动、平均换手、最大回撤和主要行业/风格暴露。汇总结果显示，Tucker 仍在 SZ50、ZZ500 和医药行业样本中兼具较好的排序和收益表现，但在 HS300 上最优累计收益与 Sharpe 来自 Tucker、最优 Rank IC 则来自 CP；在全 A 样本中收益与 Sharpe 最优为 PCA；在大市值与小市值样本中，Rank IC 和收益最优也再次发生分离。这说明模型优劣并不是单一指标可以决定的，必须同时观察排序相关性、组合收益、换手成本、超额表现和暴露集中度。

将本文结果与前述文献放在同一视角下观察，可以更清楚地说明本研究的补充意义。黄宏运与刘宇轩的研究表明，多因子模型即使在传统回归框架下，也能够通过因子筛选与宏观周期修正获得可观的选股效果；本文的结果并不否认这一点，而是进一步说明，当研究对象由“少量已筛选因子”扩展到更高维的因子集合时，问题的核心会从“因子是否有效”转向“高维信息如何被稳定压缩并用于排序”。与此同时，岳欣与曾亚丽关于张量化金融数据的研究强调了时间结构和高阶表示的重要性，而本文在样本边界扩展、长窗口年度复核以及时间状态切换结果中的发现，则进一步表明这种高阶结构优势不仅体现在波动分析或时序预测上，也可以体现在股票因子排序与模式识别任务中。换言之，本文希望在多因子研究与张量表示研究之间建立更直接的实证联系。

\section{扩展特征对照分析}

为了检验更宽的信息边界是否会改变模型表现，本文进一步构造了扩展因子面板，并在 HS300、SZ50 与 ZZ500 三个样本池上分别运行基准方案与扩展方案。扩展方案在原有四个主线因子之外，引入了市场状态代理、外部宏观变量、财务 PIT 指标、分红变量，以及业绩快报、业绩预告和公告事件字典等信息。就具体字段而言，这些新增输入覆盖了短中期收益动量、波动率、回撤、成交额变化、政策利率、价格与货币供给指标，以及盈利能力、分红行为和公告事件强度等多个维度。对照结果表明，扩展输入的影响并不呈现单边改善，而是具有明显的样本池差异。

从已生成的对照结果看，HS300 上新增市场代理变量、外部宏观、分红和公告事件字典后，三类模型的 Rank IC 全部转负，CP、Tucker 与 PCA 分别回落到 $-0.2851$、$-0.1417$ 和 $-0.1073$，组合累计收益也同步转弱，其中 CP、Tucker 与 PCA 分别为 $-1.88\%$、$-3.12\%$ 和 $-2.56\%$，说明扩展输入与大盘蓝筹样本之间仍存在显著相容性差异。SZ50 上，扩展输入并未改善排序表现，CP、Tucker 与 PCA 的 Rank IC 分别为 $-0.0345$、$-0.0509$ 与 $-0.1223$，累计收益也分别降至 $-3.07\%$、$-4.12\%$ 和 $-3.49\%$，表明“输入扩展”并没有对所有模型产生一致收益。ZZ500 上，扩展输入的分化更明显：CP 的 Rank IC 升至 0.0342，累计收益转为 1.18\%，而 PCA 与 Tucker 的 Rank IC 仍为负，累计收益分别为 $-3.47\%$ 和 $-3.39\%$。因此，扩展输入的作用更像是重新塑造不同模型在不同样本池中的信息利用方式，而不是简单地把所有结果整体抬升。

\begin{figure}[!t]
\centering
\begin{tikzpicture}
\begin{groupplot}[
  group style={group size=3 by 1, horizontal sep=1.0cm},
  width=0.29\textwidth,
  height=5.8cm,
  ybar,
  ymin=-0.20,
  ymax=0.12,
  symbolic x coords={CP,Tucker,PCA},
  xtick=data,
  enlarge y limits={upper=0.10,lower=0.08},
  enlarge x limits=0.18,
  point meta=y,
  nodes near coords={\pgfmathprintnumber[fixed,precision=2]{\pgfplotspointmeta}},
  xticklabel style={font=\tiny},
  ylabel style={font=\small},
  title style={font=\small},
  legend style={font=\tiny, draw=none, at={(0.5,1.08)}, anchor=south},
  nodes near coords,
  every node near coord/.append style={font=\tiny},
]
\nextgroupplot[title={HS300},ylabel={Rank IC}]
\addplot[fill=blue!55] coordinates {(CP,-0.0278) (Tucker,-0.0468) (PCA,-0.0530)};
\addplot[fill=orange!70] coordinates {(CP,-0.2851) (Tucker,-0.1417) (PCA,-0.1073)};
\legend{基准方案,扩展方案}
\nextgroupplot[title={SZ50}]
\addplot[fill=blue!55] coordinates {(CP,0.0198) (Tucker,0.0468) (PCA,-0.0224)};
\addplot[fill=orange!70] coordinates {(CP,-0.0345) (Tucker,-0.0509) (PCA,-0.1223)};
\nextgroupplot[title={ZZ500}]
\addplot[fill=blue!55] coordinates {(CP,-0.1082) (Tucker,0.0209) (PCA,-0.0442)};
\addplot[fill=orange!70] coordinates {(CP,0.0342) (Tucker,-0.0138) (PCA,-0.0189)};
\end{groupplot}
\end{tikzpicture}
\caption{三组正式样本池在基准方案与扩展方案下的 Rank IC 对比}
\label{fig:extended_rankic}
\end{figure}

\begin{figure}[!t]
\centering
\begin{tikzpicture}
\begin{groupplot}[
  group style={group size=3 by 1, horizontal sep=1.0cm},
  width=0.29\textwidth,
  height=5.8cm,
  ybar,
  ymin=-10.5,
  ymax=5.0,
  symbolic x coords={CP,Tucker,PCA},
  xtick=data,
  enlarge y limits={upper=0.10,lower=0.06},
  enlarge x limits=0.18,
  point meta=y,
  nodes near coords={\pgfmathprintnumber[fixed,precision=2]{\pgfplotspointmeta}},
  xticklabel style={font=\tiny},
  ylabel style={font=\small},
  title style={font=\small},
  legend style={font=\tiny, draw=none, at={(0.5,1.08)}, anchor=south},
  nodes near coords,
  every node near coord/.append style={font=\tiny},
]
\nextgroupplot[title={HS300},ylabel={累计收益(\%)}]
\addplot[fill=blue!55] coordinates {(CP,0.53) (Tucker,1.29) (PCA,0.92)};
\addplot[fill=orange!70] coordinates {(CP,-1.88) (Tucker,-3.12) (PCA,-2.56)};
\legend{基准方案,扩展方案}
\nextgroupplot[title={SZ50}]
\addplot[fill=blue!55] coordinates {(CP,-2.27) (Tucker,0.04) (PCA,-0.56)};
\addplot[fill=orange!70] coordinates {(CP,-3.07) (Tucker,-4.12) (PCA,-3.49)};
\nextgroupplot[title={ZZ500}]
\addplot[fill=blue!55] coordinates {(CP,-7.76) (Tucker,-1.63) (PCA,-1.82)};
\addplot[fill=orange!70] coordinates {(CP,1.18) (Tucker,-3.39) (PCA,-3.47)};
\end{groupplot}
\end{tikzpicture}
\caption{三组正式样本池在基准方案与扩展方案下的组合累计收益对比}
\label{fig:extended_return}
\end{figure}

图~\ref{fig:extended_rankic} 与图~\ref{fig:extended_return} 共同说明，扩展输入并没有形成“统一改善”或“统一恶化”的单边结论，而是显著增强了模型与样本池之间的交互异质性。这表明，更宽口径的信息确实会改变潜在结构，但这种改变是否能够转化为模型质量提升，仍取决于样本池风格、模型结构以及特征筛选方式。因此，本文将这组结果视为输入边界已经拓宽、但有效字段仍需继续筛选的阶段性证据。

\section{样本边界扩展探索}

为了进一步检验结论是否依赖于单一指数样本池，本文补充了样本边界扩展实验，包括全 A 活跃股票样本、三个行业分层样本以及大中小市值分层样本。与 HS300、SZ50、ZZ500 的基准样本相比，这批扩展实验更强调“不同样本边界下模型行为是否稳定”，首先在短窗口口径下完成比较，并进一步基于长窗口因子面板对 2015--2026 全部年度做年度复核，重点观察 Rank IC、滚动稳定性和组合累计收益的相对变化。

从全 A 样本结果看，Tucker 依然是三类方法中排序最稳健的一类：其 Rank IC 为 0.0279，明显优于 CP 的 $-0.0978$ 和 PCA 的 $-0.0490$；同时，其累计收益为 $-2.11\%$，虽略低于 PCA 的 $-1.82\%$，但仍显著好于 CP 的 $-5.52\%$。这说明当样本边界从指数成分池扩展到更广泛的全 A 活跃股票后，三维结构保持带来的优势并没有完全消失，而是更多体现为排序稳健性而非单一收益最优。

行业分层结果则表现出更强的结构差异。电子行业样本（C39）中，三类模型的 Rank IC 全部为负，但 CP 的 $-0.0164$ 仍明显好于 Tucker 的 $-0.1196$ 与 PCA 的 $-0.1292$，说明高波动成长行业中的截面信号更容易被特定模型结构扭曲。医药行业样本（C27）中，Tucker 的 Rank IC 达到 0.1476，累计收益达到 16.46\%，明显好于 CP 和 PCA，说明在盈利质量和事件驱动更突出的行业里，三维结构压缩与排序的协同效应更强。专用设备行业样本（C35）中，CP 的 Rank IC 最高，为 0.0520，而 Tucker 虽仅取得 $-0.0496$ 的 Rank IC，却仍保留了 0.17\% 的微弱正收益，说明不同行业边界下“最优模型”并不固定，样本风格会明显影响张量分解与二维压缩方法的相对优势。

市值分层结果进一步说明样本边界会改变排序与组合表现的对应关系。小市值样本中，CP 的 Rank IC 达到 0.1399，但组合累计收益仅为 $-1.06\%$，说明在高噪声场景下，排序相关性并不必然转化为可配置收益。中市值样本中，PCA 的 Rank IC 为 0.1224，但累计收益最优反而来自 CP 的 1.46\%，而 Tucker 仅取得 $-0.43\%$ 的轻微负收益，说明中盘样本更容易出现“排序最优”和“收益最优”分离。大市值样本中，CP 的 Rank IC 最高，为 0.1989，但三类模型组合收益均未显著转正，这再次表明“排序指标优于同业”与“组合净值最终为正”之间仍存在显著缺口。

在全年长窗口复核中，这种“最优模型随边界和年份切换”的现象依然存在。全 A 样本在 2015 窗口中呈现“Rank IC、收益与 Sharpe 最优均为 CP”的结构；2019 年三类模型的累计收益全部转正，且最优由 PCA 占据；2021 年则出现 Tucker 同时在 Rank IC、收益与 Sharpe 上占优的窗口；2026 年的短年度窗口又转为“PCA 在 Rank IC 上领先、Tucker 在收益上领先、CP 在 Sharpe 上领先”的分离格局。行业与市值分层同样没有表现出跨年份恒定不变的最优模型，说明样本边界扩展不仅改变模型排序，也改变不同市场阶段下的优势模型分布。

\begin{figure}[!t]
\centering
\begin{tikzpicture}
\begin{axis}[
  width=0.82\textwidth,
  height=7.0cm,
  ymin=0.7,
  ymax=2.4,
  symbolic x coords={2015,2019,2021,2026},
  xtick=data,
  xlabel={代表年份},
  ylabel={成本后累计净值},
  legend style={font=\small, draw=none, at={(0.5,1.05)}, anchor=south, legend columns=3},
  tick label style={font=\small},
  label style={font=\small},
  title style={font=\small},
  mark size=2.2pt,
]
\addplot+[mark=square*, blue, thick] coordinates {(2015,0.9871) (2019,1.4637) (2021,1.7381) (2026,0.9080)};
\addplot+[mark=triangle*, orange, thick] coordinates {(2015,0.8530) (2019,1.7245) (2021,1.9464) (2026,0.8600)};
\addplot+[mark=*, teal!70!black, thick] coordinates {(2015,0.6749) (2019,1.4228) (2021,2.1971) (2026,0.9194)};
\legend{CP,PCA,Tucker}
\end{axis}
\end{tikzpicture}
\caption{全 A 长窗口代表年份的成本后累计净值对比}
\label{fig:all_a_long_window_cost}
\end{figure}

图~\ref{fig:all_a_long_window_cost} 进一步给出了 2015、2019、2021 与 2026 四个代表年份下全 A 长窗口运行的成本后累计净值。该图最直接地展示了一个事实：模型优劣并不会在所有市场阶段保持不变。2015 年只有 CP 尚能把成本后净值维持在 1 附近，而 Tucker 明显落后；2019 年三类模型全部转正，但领先者改为 PCA；2021 年则出现 Tucker 明显跑赢 CP 与 PCA 的结构；到 2026 年短窗口时，三类模型又重新回落到 1 以下。换言之，长窗口复核真正补充的不是“某一模型永远最好”，而是“不同阶段的优势模型会发生切换”，这恰恰是样本边界扩展实验最重要的稳健性信息。

从暴露角度看，行业分层样本的主要行业暴露自然集中在单一行业，而全 A 与指数样本更能暴露模型选择偏好。例如全 A 的 Tucker 组合主要暴露在医药制造、黑色金属和影视传媒等行业，并且风格暴露集中在质量因子；HS300 的 Tucker 组合则明显集中于货币金融服务，同时混合质量、价值和波动率因子。图~\ref{fig:tucker_boundary_short} 进一步给出了短窗口样本边界扩展下 Tucker 的跨边界排序与收益分化：医药行业样本在 Rank IC 和累计收益上同时领先，而电子行业、小市值和大市值样本虽然保留了某些结构信号，却没有把这些信号稳定转化为正向组合净值。这个结果进一步提醒，模式发现不能只看潜在聚类是否稳定，也要检查组合层是否因为样本边界变化而集中到少数行业或单一风格。

\begin{figure}[!t]
\centering
\begin{tikzpicture}
\begin{groupplot}[
  group style={group size=2 by 1, horizontal sep=2.2cm},
  width=0.36\textwidth,
  height=6.6cm,
  xbar,
  enlarge y limits=0.18,
  yticklabel style={font=\scriptsize, text width=2.0cm, align=right},
  tick label style={font=\small},
  label style={font=\small},
  title style={font=\small},
  point meta=x,
  nodes near coords={\pgfmathprintnumber[fixed,precision=2]{\pgfplotspointmeta}},
  nodes near coords,
  every node near coord/.append style={font=\tiny},
]
\nextgroupplot[
  title={2026 长窗口 Tucker Rank IC},
  xmin=-0.16,
  xmax=0.14,
  symbolic y coords={电子C39,小市值,全市场,专用设备C35,大市值,医药C27,中市值},
  ytick=data,
  xlabel={Rank IC}
]
\addplot[fill=teal!65] coordinates {
(-0.1215,电子C39)
(-0.0877,小市值)
(-0.0708,全市场)
(-0.0462,专用设备C35)
(0.0102,大市值)
(0.0161,医药C27)
(0.0814,中市值)
};
\nextgroupplot[
  title={2026 长窗口 Tucker 期末净值},
  xmin=0.68,
  xmax=1.30,
  symbolic y coords={小市值,电子C39,中市值,专用设备C35,全市场,大市值,医药C27},
  ytick=data,
  xlabel={期末净值}
]
\addplot[fill=purple!45] coordinates {
(0.7940,小市值)
(0.8156,电子C39)
(0.8745,中市值)
(0.8807,专用设备C35)
(0.9196,全市场)
(0.9260,大市值)
(1.2818,医药C27)
};
\end{groupplot}
\end{tikzpicture}
\caption{2026 年长窗口结果中 Tucker 的跨边界表现}
\label{fig:tucker_boundary_long2026}
\end{figure}

\section{稳健性与局限性分析}

在滚动稳定性方面，HS300、SZ50 与 ZZ500 上的 Tucker 分别达到 0.9898、0.9923 和 0.9683，均显著高于 CP，也与 PCA 基本处于同一量级。稳定性的意义在于，模型在不同滚动窗口上学到的潜在结构具有较强一致性，这使得研究者能够更放心地解释因子重要性和时间状态变化，而不会因为模型在不同窗口下频繁漂移而失去可解释基础。对研究型实证而言，稳定性往往比单次高分更重要，因为它决定了结论是否具备可重复性。

\begin{figure}[htbp]
\centering
\begin{tikzpicture}
\begin{groupplot}[
  group style={group size=3 by 1, horizontal sep=1.3cm},
  width=0.28\textwidth,
  height=6.0cm,
  ybar,
  ymin=0.65,
  ymax=1.02,
  symbolic x coords={CP,Tucker,PCA},
  xtick=data,
  nodes near coords,
  nodes near coords align={vertical},
  every node near coord/.append style={font=\tiny},
  tick label style={font=\small},
  ylabel style={font=\small},
  title style={font=\small}
]
\nextgroupplot[title={HS300},ylabel={滚动稳定性}]
\addplot[fill=blue!55] coordinates {(CP,0.6974) (Tucker,0.9898) (PCA,0.9946)};
\nextgroupplot[title={SZ50}]
\addplot[fill=orange!70] coordinates {(CP,0.5320) (Tucker,0.9923) (PCA,0.9956)};
\nextgroupplot[title={ZZ500}]
\addplot[fill=teal!65] coordinates {(CP,0.8384) (Tucker,0.9683) (PCA,0.9575)};
\end{groupplot}
\end{tikzpicture}
\caption{三个正式样本池上的模型滚动稳定性对比图}
\label{fig:stability}
\end{figure}

当然，本文仍存在若干需要继续完善的方面。第一，在数据输入层面，当前训练输入已经在扩展方案中纳入市场状态代理、外部利率、宏观月度指标、财务 PIT、分红以及公告事件字典特征，并形成了较清晰的披露时点与可用时点约束。但外部宏观与公告文本的覆盖仍以当前实验窗口附近为主，尚未完全扩展到更长历史与更完整正文全文，这意味着当前模型虽然已经验证了更宽输入边界可以在点时一致的前提下进入训练样本，但对于更宽口径、多来源因子体系的长期适应性仍需在源表持续扩充后继续检验。

第二，在效果验证层面，本文已经补齐了前N组合的累计净值、回撤、换手率和行业/风格暴露等基础结果，并进一步加入了分位数组、多空组合、成本后净值、交易成本和相对样本池指数的超额收益，且已通过跨样本边界汇总表与 \texttt{2015-2026} 全年长窗口组合汇总统一比较。该证据链已经能够从短窗口与全年窗口两个层面回答“若将排序结果转化为真实组合规则，会得到怎样的收益—风险画像”；当前仍需继续深化的是更细粒度的交易冲击刻画和风险归因模型，而不是再停留在“是否具备长窗口证据”这一层面。

第三，在样本边界层面，本文已经从 HS300、SZ50 与 ZZ500 三个指数样本池扩展到全 A 活跃股票、行业分层和市值分层实验，并为 2015--2026 年区间完成了按年度拆分的长窗口复核。新增结果说明，指数样本上的结论不能被机械外推：Tucker 在部分窗口和部分行业样本中表现较稳健，但电子、专用设备和不同市值分层样本的最优模型并不固定。虽然全年长窗口结果已经补齐，但当前训练输入仍属于既定研究边界下的正式比较，因此这批扩展实验更适合被理解为“现有输入范围在不同边界与市场阶段下的稳健性检验”，而不是所有候选输入都已完备后的最终全市场结论。

第四，在可视化与论证层面，本文已经补充解释方差、Rank IC、滚动稳定性、时间状态、因子热力图、股票潜在结构图、行业聚类交叉图和跨样本边界对比图，使模式发现叙事能够由图表与文字共同支撑。进一步补充的 2015--2026 各年度全 A 长窗口时间状态切换图、因子重要性热力图及年度比较图，已经把此前“长窗口图待补”的占位说明替换为真实证据。当前仍需继续完善的是把这些长窗口结果与更宽输入边界实验联动，而不是继续停留在主线四因子输入上。

针对上述不足，本文给出如下完善计划。其一，在数据层面，继续按照“先明确可用时点，再逐步扩展特征输入”的顺序推进：已完成市场状态代理、外部宏观、财务 PIT、分红和公告事件字典的第一阶段披露时点校验，后续重点转向更长历史、更多宏观品种以及公告正文全文。其二，在评估层面，现有分析已经覆盖前N组合、分组收益、多空、成本后收益、超额收益、回撤、换手、波动、Sharpe 和暴露，后续重点转向字段筛选、长窗口风险归因与交易冲击刻画。其三，在样本层面，保留指数样本池实验作为稳定基线，同时使用全 A 股、行业分层和不同市值分层实验比较不同样本边界下潜在结构、最优秩和稳定性指标的差异。其四，在结果表达层面，继续扩展长窗口图形体系，并进一步增强图表与文字论证之间的一致性。

\chapter{结论与展望}

\section{研究结论}

本文以“股票—因子—时间”三维张量为统一研究对象，围绕股票因子降维与模式发现问题，构建了正式数据链、张量建模框架和结构化实验输出体系。基于 HS300、SZ50 与 ZZ500 三个正式样本池在短窗口输出口径下的实验结果，本文得到以下主要结论。

第一，股票因子问题天然具有多模态结构，使用三阶张量表示比单纯二维矩阵表示更符合数据本质。张量表示不仅保留了股票截面、因子集合和时间演化的耦合关系，也为后续模式发现提供了天然结构基础。

第二，从数值重构角度看，PCA 作为二维压缩基线依然具有显著优势；但从结构保持与模式解释角度看，Tucker 分解在三个样本池上表现出更平衡的重构质量、排序能力与长期稳定性。尤其在 Rank IC 和滚动稳定性指标上，Tucker 相较于 CP 更稳定，也比单纯依赖 PCA 更能支撑结构解释与排序分析。

第三，张量分解不仅可以完成因子降维，还能够输出股票相似结构、因子共振关系和时间状态变化等模式信息。实验结果显示，质量、价值、动量和波动率等因子在潜在空间中承担了较重要的解释角色，不同时间阶段也存在显著的结构重排现象，说明张量方法适合从“模式发现”而不仅仅是“数值压缩”的视角理解股票因子系统。

\section{创新点与不足}

本文的创新主要体现在三个方面。其一，将股票因子降维、模式发现与 A 股样本实证放在统一框架下讨论，尽量避免算法分析、数据处理与结论表达彼此割裂。其二，在毕业论文层面同时比较 CP、Tucker 与 PCA 三类方法，不仅比较重构误差，还比较 Rank IC 与滚动稳定性，使结论更贴近金融研究需求。其三，在指数样本、行业样本、市值样本以及长窗口复核之间建立了相互对应的比较关系，使论文中的图表、结果表和文字论证能够形成较为一致的证据链条。

从研究内容看，本文并未停留在“少量因子是否有效”的传统讨论上，而是把关注点前移到高维因子集合的表示学习与结构压缩问题。与此同时，本文也没有把张量分解仅仅用作数值压缩工具，而是进一步讨论其在横截面排序、样本边界扩展、行业聚类和时间状态识别中的解释作用。通过这种安排，论文希望把多因子研究中的选股问题与张量方法中的高阶结构分析更自然地结合起来。

不足之处主要包括：虽然全 A、行业和市值分层实验已经补充到 2015--2026 全部年度的长窗口复核，并形成了相应的跨年度汇总表与比较图，但外部宏观、分红、重大事项和公告文本等补充信息仍以当前实验窗口附近的数据覆盖为主，尚未扩展到更长历史与更完整文本层面；当前实验已经形成前N组合净值、分位数组、多空、成本后净值、基准超额收益和跨样本边界统一汇总，但交易冲击和风险归因仍有继续细化空间；模式发现图表已经从基础指标扩展到股票潜在结构、行业聚类关系和全年长窗口结果，但其中仍有进一步压缩叙述、增强图文配合的空间。

\section{后续工作展望}

后续工作可以从三个方向继续推进。第一，在数据层继续扩展已接入的外部宏观、分红、重大事项和公告文本源表，使其覆盖更长历史、更完整公告正文以及更细的事件分类，同时保持 PIT 安全。第二，在方法层继续研究秩自动选择、稀疏约束、非负张量分解和动态张量分解，以提升模型解释性与适应性。第三，在应用层把当前结构化实验输出延伸到长窗口组合复核、回撤控制、风险归因和风格轮动监测，从而形成更完整的量化研究闭环。

\chapter*{参考文献}
\addcontentsline{toc}{chapter}{参考文献}
\sloppy
\begin{enumerate}[label={[\arabic*]}]
  \item Hitchcock F L. The Expression of a Tensor or a Polyadic as a Sum of Products[J]. Journal of Mathematical Physics, 1927, 6(1): 164-189.
  \item Tucker L R. Some Mathematical Notes on Three-Mode Factor Analysis[J]. Psychometrika, 1966, 31: 279-311.
  \item Harshman R A. Foundations of the PARAFAC Procedure: Models and Conditions for an Explanatory Multi-modal Factor Analysis[J]. UCLA Working Papers in Phonetics, 1970, 16: 1-84.
  \item Kolda T G, Bader B W. Tensor Decompositions and Applications[J]. SIAM Review, 2009, 51(3): 455-500.
  \item 胡胜龙. 张量分解的唯一性[J]. 运筹学学报(中英文), 2025, 29(03): 34-60.
  \item 夏虹, 张雅倩, 靳晓东, 陈彦萍, 高聪, 王忠民. 基于张量的方法及应用综述[J]. 计算机技术与发展, 2022, 32(06): 1-8.
  \item Rabanser S, Shchur O, Günnemann S. Introduction to Tensor Decompositions and their Applications in Machine Learning[EB/OL]. arXiv:1711.10781, 2017.
  \item 岳欣. 基于张量分解的股票波动分析[D]. 西南财经大学, 2022.
  \item 马少沛, 孙庆慧, 武雅萱, 等. 大数据下张量充分降维方法及其应用研究[J]. 统计研究, 2021, 38(02): 114-134.
  \item 曾亚丽. 基于张量模型的金融时间序列预测研究[D]. 武汉理工大学, 2017.
  \item Brandi G, Gramatica R, Di Matteo T. Unveil Stock Correlation via a New Tensor-based Decomposition Method[EB/OL]. arXiv:1911.06126, 2020.
  \item Spelta A. Financial Market Predictability with Tensor Decomposition and Links Forecast[J]. Applied Network Science, 2017, 2: 7.
  \item 黄宏运, 王梅, 朱家明. 基于多元回归分析的多因子选股模型[J]. 通化师范学院学报, 2016, 37(08): 44-46.
  \item 刘宇轩, 金伟泽, 袁亮. 多因子量化选股模型优化与实证研究——引入金融周期指标的分析[J]. 价格理论与实践, 2022(04): 141-145.
  \item Lettau M. High-Dimensional Factor Models with an Application to Mutual Fund Characteristics[R]. NBER Working Paper No.29833, 2022.
  \item Han Y F, Chen R, Zhang C H. Rank Determination in Tensor Factor Model[EB/OL].\\
  arXiv:2011.07131, 2022.
  \item Wang D, Zheng Y, Lian H, Li G D. High-dimensional Vector Autoregressive Time Series Modeling via Tensor Decomposition[EB/OL]. arXiv:1909.06624, 2020.
  \item 李博, 孟志青, 朱爱花. 时态支持向量机模型在股票操纵模式发现上的研究[J]. 系统科学与数学, 2023, 43(02): 356-378.
  \item 徐常青, 祁力群. 张量CP分解、半正定张量和范德蒙张量[J]. 苏州科技学院学报(自然科学版), 2016, 33(02): 1-7+82.
\end{enumerate}
\fussy

\appendix
\chapter{实验参数与补充说明}

本文正式实验以 HS300、SZ50 与 ZZ500 为样本池。正文首先讨论 2026 年 3 月 2 日至 2026 年 3 月 30 日窗口下的短窗口结果，模型层固定比较 CP、Tucker 与 PCA 三类方法，评价指标包括 MSE、解释方差、Rank IC、IR 与滚动稳定性。实验结果通过统一的结构化评价流程整理，确保论文中的图表、表格和文本结论具有一致证据来源。

\chapter{关键程序代码附录}

为便于说明正文中的关键约束与核心算法，本文选取与主要结论直接相关的程序片段作为附录。附录代码仅保留能够支撑方法解释的关键逻辑，其目的在于说明研究设计如何在实现层面得到落实，而不替代正文中的理论推导与实验分析。

\section{CP 分解的交替最小二乘更新}

下列代码片段对应 CP 分解训练中的核心迭代部分。其基本思路是：在固定另外两个模态因子矩阵后，分别对股票模态、因子模态和时间模态做交替更新，并在每轮迭代后进行列归一化，从而得到可比较的潜在成分权重。

\begin{lstlisting}[language=Python]
for iteration in range(1, max_iter + 1):
    # 固定因子模态与时间模态后更新股票模态载荷。
    gram = (c_mat.T @ c_mat) * (b_mat.T @ b_mat) + identity
    a_mat = (unfold(tensor, 0) @ khatri_rao(c_mat, b_mat)) @ np.linalg.pinv(gram)

    # 固定股票模态与时间模态后更新因子模态载荷。
    gram = (c_mat.T @ c_mat) * (a_mat.T @ a_mat) + identity
    b_mat = (unfold(tensor, 1) @ khatri_rao(c_mat, a_mat)) @ np.linalg.pinv(gram)

    # 固定股票模态与因子模态后更新时间模态载荷。
    gram = (b_mat.T @ b_mat) * (a_mat.T @ a_mat) + identity
    c_mat = (unfold(tensor, 2) @ khatri_rao(b_mat, a_mat)) @ np.linalg.pinv(gram)

    weights = np.ones(rank, dtype=float)
    for factor_matrix in (a_mat, b_mat, c_mat):
        # 列归一化使成分尺度集中到标量权重中，便于比较。
        norms = np.linalg.norm(factor_matrix, axis=0)
        norms[norms == 0] = 1.0
        factor_matrix /= norms
        weights *= norms
\end{lstlisting}

\section{Tucker 分解的多模态交替更新}

下列代码片段体现了 Tucker 分解相对 CP 分解更灵活的地方：不同模态可以设置不同秩，并通过对各模态展开矩阵做奇异向量更新，交替优化股票、因子与时间三个方向的低维表示。核张量 \texttt{core} 用于刻画这些低维表示之间的高阶交互。

\begin{lstlisting}[language=Python]
u_stock = _top_singular_vectors(unfold(tensor, 0), stock_rank)
u_factor = _top_singular_vectors(unfold(tensor, 1), factor_rank)
u_time = _top_singular_vectors(unfold(tensor, 2), time_rank)

for iteration in range(1, max_iter + 1):
    # 在固定因子模态与时间模态后更新股票模态。
    projected = _mode_product(_mode_product(tensor, u_factor.T, 1), u_time.T, 2)
    u_stock = _top_singular_vectors(unfold(projected, 0), stock_rank)

    # 在固定股票模态与时间模态后更新因子模态。
    projected = _mode_product(_mode_product(tensor, u_stock.T, 0), u_time.T, 2)
    u_factor = _top_singular_vectors(unfold(projected, 1), factor_rank)

    # 在固定股票模态与因子模态后更新时间模态。
    projected = _mode_product(_mode_product(tensor, u_stock.T, 0), u_factor.T, 1)
    u_time = _top_singular_vectors(unfold(projected, 2), time_rank)

    # 根据三组模态载荷反推出核张量，保留高阶交互信息。
    core = _mode_product(
        _mode_product(_mode_product(tensor, u_stock.T, 0), u_factor.T, 1),
        u_time.T,
        2,
    )
\end{lstlisting}

\section{PCA 基线的张量展开与重构}

为了与张量方法形成可比的二维压缩基线，本文并未直接在原始表格上做 PCA，而是先把三阶张量沿时间和因子方向展开，再对展开矩阵进行奇异值分解。训练完成后，再把低维结果重构回股票—因子—时间三维张量，以保证评价口径与 CP、Tucker 保持一致。

\begin{lstlisting}[language=Python]
matrix = np.transpose(tensor, (0, 2, 1)).reshape(
    stock_count * time_count,
    factor_count,
)
mean_vector = matrix.mean(axis=0, keepdims=True)
centered = matrix - mean_vector
left, singular_values, right_t = np.linalg.svd(
    centered,
    full_matrices=False,
)

# 以前 rank 个奇异值和奇异向量构造二维低维表示。
scores = left[:, :rank] * singular_values[:rank]
components = right_t[:rank, :].T

# 将二维重构结果恢复到三维张量形状，以保持评价口径一致。
reconstruction_matrix = scores @ components.T + mean_vector
reconstruction = np.transpose(
    reconstruction_matrix.reshape(stock_count, time_count, factor_count),
    (0, 2, 1),
)
\end{lstlisting}

\section{扩展特征的可用时点与时效计算}

下面的代码片段展示了扩展因子面板中“特征年龄”与事件时效衰减的处理方式。其关键点在于：所有时效性特征均从 \texttt{available\_date} 而非 \texttt{pub\_date} 起算，从而保证披露日落在非交易日或收盘后时，模型不会提前获得未来信息。

\begin{lstlisting}[language=Python]
def _snapshot_age_days(snapshot, trade_date: str) -> int:
    if snapshot is None:
        return 0
    # 以最早可交易时点作为年龄起点，保证时点安全。
    effective_date = snapshot.available_date or snapshot.pub_date
    return _days_since(effective_date, trade_date)


def _event_age_decay_score(
    performance_snapshot: PerformanceExpressSnapshot | None,
    forecast_snapshot: ForecastSnapshot | None,
    trade_date: str,
) -> float:
    score = 0.0
    if performance_snapshot is not None:
        # 快报越新，衰减项越大。
        score += 1.0 / (
            1.0 + _snapshot_age_days(performance_snapshot, trade_date)
        )
    if forecast_snapshot is not None:
        # 预告按相同规则计入时间衰减。
        score += 1.0 / (
            1.0 + _snapshot_age_days(forecast_snapshot, trade_date)
        )
    return score
\end{lstlisting}

\section{扩展源表的公告日映射与去重}

下段代码体现了扩展源表构建时的两个关键原则：其一，窗口前事件保持原始公告日，避免历史信息被错误压到首个交易日；其二，公告表按“股票代码—公告类型—公告日—标题—链接”去重，防止同一公告被重复统计到事件特征中。

\begin{lstlisting}[language=Python]
def _source_available_date(pub_date: str, trade_dates: list[str]) -> str:
    # 若公告早于样本窗口的首个交易日，则保留原始日期。
    if trade_dates and pub_date < trade_dates[0]:
        return pub_date
    return _next_trade_date_after(pub_date, trade_dates) or ""


def _notice_identity(
    row: dict[str, object],
) -> tuple[str, str, str, str, str]:
    # 用于识别重复公告，避免同一事件被重复统计到。
    return (
        str(row.get("stock_code", "")),
        str(row.get("notice_type", "")),
        str(row.get("pub_date", "")),
        str(row.get("title_text", "")),
        str(row.get("url", "")),
    )
\end{lstlisting}
